\subsection{メモリ}

時々C++プログラマから\q{スタックにメモリを割り当てる}や\q{\gls{heap}にメモリを割り当てる}という言葉を聞くことがあります。

オブジェクトを\IT{スタック上で}割り当てる：

\begin{lstlisting}[style=customc]
void f()
{
	...

	Class o=Class(...);

	...
};
\end{lstlisting}

オブジェクト（または構造体）のメモリは、単純な\ac{SP}のシフトを使用してスタックに割り当てられます。
メモリは関数の終了時、またはより正確には\IT{スコープ}の終了時に解放されます---\ac{SP}が（関数の開始時と同じ）状態に戻り、\IT{Class}のデストラクタが呼び出されます。
同様に、Cで割り当てられた構造体のメモリは関数の終了時に解放されます。

\IT{\gls{heap}上で}オブジェクトを割り当てる：

\begin{lstlisting}[style=customc]
void f1()
{
	...

	Class *o=new Class(...);

	...
};

void f2()
{
	...

	delete o;

	...
};
\end{lstlisting}

\myindex{\CStandardLibrary!malloc()}
\myindex{\CStandardLibrary!free()}
これは\IT{malloc()}呼び出しを使用して構造体のメモリを割り当てるのと同じです。
実際、C++の\IT{new}は\IT{malloc()}のラッパーであり、\IT{delete}は\IT{free()}のラッパーです。
メモリブロックが\gls{heap}に割り当てられているため、\IT{delete}を使用して明示的に解放する必要があります。
クラスのデストラクタはその直前に自動的に呼び出されます。

どちらの方法が良いでしょうか？
\IT{スタック上での}割り当ては非常に高速で、現在の関数でのみ使用される小さな短命のオブジェクトに適しています。

\IT{ヒープ上での}割り当ては遅く、多くの関数で使用される長命のオブジェクトに適しています。
また、\gls{heap}に割り当てられたオブジェクトは、明示的に解放する必要があるため、メモリリークが発生しやすくなりますが、忘れてしまうことがあります。

いずれにしても、これは好みの問題です。